\chapter{Wendelin Werner}\label{ch2006wendelinwerner:wendelin-werner-fields-medal-2006-a-comprehensive-report-on-his-work-and-impact}

\begin{center}\rule{0.5\linewidth}{0.5pt}\end{center}

\textbf{Prepared:} July 2026 \textbf{Subject:} Wendelin Werner (b. September 23, 1968, Cologne, West Germany) \textbf{Position:} Professor of Mathematics, ETH Zürich (since 2006) \textbf{Doctoral Advisor:} Jean-François Le Gall (Université Paris VI) \textbf{Fields Medal Citation:} \emph{``For his contributions to the development of stochastic Loewner evolution, the geometry of two-dimensional Brownian motion, and conformal field theory.''}

\begin{center}\rule{0.5\linewidth}{0.5pt}\end{center}

\begin{prereqs} Probability theory (Brownian motion, Markov processes), basic complex analysis (conformal maps, Riemann mapping theorem), introductory statistical mechanics. \end{prereqs}

\begin{objectives} Define the Schramm--Loewner evolution and its parameter $\kappa$ \cite{werner2004}; prove the exact Hausdorff dimension of SLE curves; understand the conformal restriction property and the Brownian loop soup. \end{objectives}

\section{Problem History}\label{ch2006wendelinwerner:problem-history}

A central problem in probability theory and statistical mechanics is to identify the scaling limits of discrete random models at criticality. Two-dimensional models such as critical percolation on the triangular lattice, the Ising model at its critical temperature, and the self-avoiding walk all exhibit scale invariance at criticality. Physicists long conjectured that these models should have \emph{conformally invariant} scaling limits, meaning that the limiting objects are unchanged under conformal (angle-preserving) transformations of the domain.

The difficulty lay in constructing these limits rigorously. Conformal invariance cannot be imposed on discrete models directly; it must emerge in the scaling limit. Without a mathematical framework for conformally invariant random curves, progress was limited to heuristic arguments borrowed from conformal field theory (CFT).

In 1999, Oded Schramm introduced a one-parameter family of random curves in the plane by driving the classical Loewner equation with a Brownian motion. This construction, now called the \textbf{Schramm--Loewner evolution} (SLE), provided the first explicit candidates for conformally invariant random curves. The parameter $\kappa > 0$ controls the law of the curve, and different values of $\kappa$ correspond to the scaling limits of different statistical mechanics models.

The problem was to develop the mathematical theory of SLE, to compute its geometric and probabilistic properties exactly, and to use SLE to identify and analyze the scaling limits of concrete models. The theory needed to be developed from scratch: Schramm's original paper was only three pages long, and the mathematical toolkit for studying these random curves did not yet exist.

\section{How Was It Solved}\label{ch2006wendelinwerner:how-was-it-solved}

\textbf{The Loewner Equation and SLE.}

Charles Loewner (1923) showed that any simple curve $\gamma$ growing from a boundary point of the upper half-plane $\mathbb{H}$ determines a family of conformal maps $g_t: \mathbb{H} \setminus \gamma([0,t]) \to \mathbb{H}$ satisfying the differential equation

\[
\frac{\partial}{\partial t} g_t(z) = \frac{2}{g_t(z) - U_t}, \qquad g_0(z) = z,
\]

where $U_t = g_t(\gamma(t))$ is the \emph{driving function} on $\partial\mathbb{H}$. Schramm's insight was to set $U_t = \sqrt{\kappa}\, B_t$, where $B_t$ is a standard Brownian motion. The resulting random curve $\gamma_\kappa$ is the \textbf{SLE$_\kappa$}. The value of $\kappa$ determines the geometry of the curve:

\begin{itemize}
  \item $0 < \kappa \leq 4$: $\gamma_\kappa$ is a simple (non-self-intersecting) curve.
  \item $4 < \kappa < 8$: $\gamma_\kappa$ is self-intersecting but does not fill the plane.
  \item $\kappa \geq 8$: $\gamma_\kappa$ is space-filling.
\end{itemize}

The dichotomy at $\kappa = 4$ arises because the driving Brownian motion $B_t$ changes sign infinitely often on every time interval; when $\sqrt{\kappa} > 4$, the curve can ``reach back'' and hit itself.

\textbf{Hausdorff Dimension of SLE$_\kappa$.}

Werner, together with Lawler and Schramm, proved that the Hausdorff dimension of the trace of SLE$_\kappa$ is

\[
\dim_H(\text{trace of SLE}_\kappa) = 1 + \frac{\kappa}{8} \qquad \text{for } \kappa \leq 8.
\]

The proof uses the relationship between the hitting probability of small sets and the Green's function of the driving Brownian motion. For $\kappa \leq 4$, the curve is simple and the dimension formula follows from the Beffara bound combined with upper and lower estimates on the probability that the curve passes near a given point. For $\kappa > 4$, the self-intersections must be controlled using the support theorem for SLE.

\textbf{Brownian Intersection Exponents.}

Werner's PhD work, continued in collaboration with Lawler and Schramm, established exact values for the \emph{intersection exponents} of planar Brownian motion. For $p$ independent Brownian motions starting from the origin, the probability that they do not intersect up to time $t$ decays as

\[
\mathbb{P}(\text{no intersection up to time } t) \asymp t^{-\xi(p)},
\]

where the exponent is given by

\[
\xi(p) = \frac{4p^2 - 1}{24}.
\]

This formula had been predicted by physicists using CFT with central charge $c = 0$; the proof uses the connection between Brownian intersection exponents and SLE partition functions, combined with the conformal martingale property of certain SLE-related quantities.

\textbf{The Conformal Restriction Property.}

Werner, with Lawler and Schramm, characterized all random connected sets $K$ in $\mathbb{H}$ satisfying the \emph{conformal restriction property}: for any hull $A$ with $K \cap A = \emptyset$, the conditional law of $K$ in $\mathbb{H} \setminus A$ is the same as the law of $K$ in $\mathbb{H}$, transported by the conformal map. They proved that the only such objects are:

\begin{enumerate}
  \item SLE$_\kappa$ for $\kappa \in (8/3, 4]$, and
  \item the \emph{Brownian loop soup} with parameter $c$ in a certain range.
\end{enumerate}

This classification provides a powerful uniqueness tool: any random curve satisfying the restriction property must be SLE$_{8/3}$, the unique simple curve with this property.

\textbf{The Brownian Loop Soup and CLE.}

Lawler and Werner introduced the \textbf{Brownian loop soup}, a Poisson point process of Brownian loops in a domain, with a conformally invariant intensity measure. Werner and Sheffield later defined the \textbf{Conformal Loop Ensemble} (CLE$_\kappa$), the collection of outermost loop boundaries of the loop soup, and proved it is the unique conformally invariant measure on non-intersecting loops satisfying a natural Markov property. The CLE provides a framework for understanding multiple non-intersecting interfaces in statistical mechanics.

\section{Impact of This Discovery}\label{ch2006wendelinwerner:impact-of-this-discovery}

The development of SLE had a transformative impact across probability theory, mathematical physics, and complex analysis.

In \textbf{probability theory}, SLE provided the first rigorous construction of conformally invariant random curves and established them as central objects of study. The exact computation of geometric invariants (Hausdorff dimension, intersection exponents) showed that conformal invariance determines these quantities uniquely, without detailed lattice calculations. The theory of SLE martingales and their connection to the Poisson boundary gave new tools for studying two-dimensional random processes.

In \textbf{statistical mechanics}, SLE identified the scaling limits of the critical Ising model interfaces (SLE$_3$), critical percolation exploration paths (SLE$_6$), loop-erased random walk (SLE$_2$), and provided the conjectured limit for the self-avoiding walk (SLE$_{8/3}$). These identifications, combined with exact computations of crossing probabilities and boundary exponents, settled questions that had been open for decades.

In \textbf{mathematical physics}, the connection between SLE and conformal field theory was made rigorous. The intersection exponents of Brownian motion correspond to the operator product expansion of the CFT with $c = 0$, and the partition functions of multiple SLE curves satisfy the Knizhnik--Zamolodchikov equations. This established a rigorous bridge between probability and the representation theory of the Virasoro algebra.

The \textbf{Brownian loop soup} and CLE provide a geometric framework for understanding conformal field theories at central charge $c \leq 1$. The loop soup is the fundamental conformally invariant measure on loops, and its cluster structure encodes the SLE curves. For central charge $c \in [0, 1]$, the loop soup parameter and the SLE parameter are related by $c = (6-\kappa)(3\kappa-8)/(2\kappa)$. The \emph{central charge} $c$ is a fundamental parameter of the conformal field theory that measures the ``degrees of freedom'' of the theory. For the Ising model, $c = 1/2$ and $\kappa = 3$; for percolation, $c = 0$ and $\kappa = 6$; for the self-avoiding walk, $c = 0$ and $\kappa = 8/3$. The formula $c = (6-\kappa)(3\kappa-8)/(2\kappa)$ is the \emph{Cardy formula} relating the loop soup parameter to the SLE parameter, and it provides a precise connection between the probabilistic (SLE) and algebraic (CFT) descriptions of the same conformally invariant object.

\section{Modern Views and Improvements}\label{ch2006wendelinwerner:modern-views-and-improvements}

The theory of SLE has been extended in several directions. The \textbf{imaginary geometry} program of Miller and Sheffield constructs multiple SLE curves as flow lines of the Gaussian free field, providing a unified framework for understanding interactions between SLE curves. The \textbf{mating of trees} bijection, developed by Duplantier, Miller, and Sheffield, connects SLE$_\kappa$ to correlated Brownian motions, yielding new results on the geometry of SLE and the Brownian map.

The \textbf{Liouville quantum gravity} (LQG) program, building on work of Duplantier and Sheffield, constructs random surfaces using the Gaussian free field. The relationship $\gamma = \sqrt{\kappa}$ connects the LQG parameter $\gamma$ to the SLE parameter $\kappa$, and recent results have identified the scaling limits of random planar maps as LQG surfaces decorated by SLE curves.

Open problems include: identifying the scaling limit of the self-avoiding walk as SLE$_{8/3}$ (proved by Duminil-Copin and Smirnov in the case of the hexagonal lattice), extending SLE to higher dimensions (where conformal invariance no longer holds), and developing a complete classification of conformally invariant measures on random curves and surfaces.

The \emph{SLE$_\kappa$ for $\kappa > 8$} produces space-filling curves that touch every point in the domain. The trace of SLE$_\kappa$ for $\kappa > 8$ is a space-filling curve whose interior is a random fractal set with Hausdorff dimension $1 + 2/\kappa$. The space-filling property means that the curve visits every point in the domain, and the interior is a random set with a complicated fractal structure. The SLE$_\kappa$ for $\kappa > 8$ is relevant for the scaling limit of the harmonic explorer and for the study of the Gaussian free field level lines. The phase transition at $\kappa = 8$ is sharp: for $\kappa < 8$ the curve is not space-filling, while for $\kappa \ge 8$ it is.

The \emph{scaling limit} of the self-avoiding walk on the hexagonal lattice was proved to be SLE$_{8/3}$ by Duminil-Copin and Smirnov (2012). The proof uses the discrete holomorphic observable for the self-avoiding walk, which is a sum over self-avoiding paths weighted by a complex weight depending on the winding angle. The observable satisfies a discrete holomorphicity condition on the hexagonal lattice (but not on other lattices), and its scaling limit is a holomorphic function whose boundary values determine the crossing probability. The result confirms a prediction of Conformal Field Theory and provides the first rigorous proof of conformal invariance for the self-avoiding walk.

The \textbf{Conformal Loop Ensemble} has become the natural framework for studying multiple interfaces. Recent work by Kemppainen and Smirnov establishes the existence of CLE as the scaling limit of the Fortuin--Kasteleyn--Renyaard cluster boundaries in random cluster models.

The \emph{Liouville quantum gravity} (LQG) surface is a random surface equipped with a conformal structure induced by the Gaussian free field. The LQG parameter $\gamma \in (0,2)$ controls the roughness of the surface, and the relationship $\gamma = \sqrt{\kappa}$ connects the LQG parameter to the SLE parameter $\kappa$. For $\gamma = \sqrt{8/3}$, the LQG surface decorated by SLE$_{8/3}$ curves describes the scaling limit of the self-avoiding walk on a random planar map. The LQG surface is defined as the ensemble of random measures $e^{\gamma h(z)} dz$ where $h$ is the Gaussian free field, and the conformal structure is induced by the GFF via the coordinate changes $e^{\gamma h(z)} dz$. The KPZ (Kardar--Parisi--Zhang) formula relates the fractal dimension of a set on an LQG surface to its dimension on the underlying Gaussian free field, and this formula is essential for the study of SLE curves on LQG surfaces.

Recent progress on the \emph{Schramm--Loewner evolution for $d$-dimensional spaces} and the \emph{massive SLE} (SLE with a mass term) continue to extend the reach of Werner's foundational contributions. The relationship between SLE and the KPZ (Kardar--Parisi--Zhang) universality class for random growth processes is another active direction, connecting Werner's work to problems in non-equilibrium statistical mechanics.

The \emph{Schramm--Loewner evolution} can be extended to settings beyond the upper half-plane. On a Riemann surface $\Sigma$ with boundary, the SLE$_\kappa$ is defined by mapping the surface to $\mathbb{H}$ using the Riemann mapping theorem and pulling back the SLE$_\kappa$ from $\mathbb{H}$. The resulting curve depends on the conformal type of $\Sigma$, and the conformal invariance of SLE ensures that the law of the curve is determined by the conformal class of the initial data. For \emph{radial SLE$_\kappa$}, the curve grows from a boundary point toward an interior point, and the driving function satisfies a different Loewner equation. The radial SLE$_\kappa$ is relevant for the scaling limit of loop-erased random walk and the uniform spanning tree, and its geometric properties differ from chordal SLE$_\kappa$ in important ways (e.g., the Hausdorff dimension of radial SLE$_\kappa$ is $1 + \kappa/8$ for $\kappa < 8$, the same as chordal).

The \emph{local time} of an SLE$_\kappa$ curve measures the amount of time the curve spends near a given point. For $\kappa > 4$, the SLE$_\kappa$ curve hits itself and can be parameterized by its self-intersections. The \emph{self-intersection local time} is a stochastic process that measures the total amount of self-intersection up to time $t$, and it satisfies a stochastic differential equation driven by the driving Brownian motion $B_t$. The local time is used to define the \emph{filled hull} of the SLE$_\kappa$ curve (the set of points separated from infinity by the curve), and the geometry of the filled hull is central to the study of SLE$_\kappa$ for $\kappa > 4$.

The \emph{Green's function} of SLE$_\kappa$ measures the probability that the curve passes near a given point. For a point $z \in \mathbb{H}$, the Green's function is $G(z) = \lim_{\varepsilon \to 0} \varepsilon^{d-2} \mathbb{P}(z \in \text{hull of SLE at capacity } t)$, where $d = 1 + \kappa/8$ is the Hausdorff dimension. The Green's function satisfies a second-order ODE (the \emph{Gambler's ruin equation}) and is explicitly computable in terms of hypergeometric functions. The Green's function is used to compute intersection exponents, crossing probabilities, and other geometric quantities for SLE$_\kappa$.

\begin{center}\rule{0.5\linewidth}{0.5pt}\end{center}

\section*{Summary}
\addcontentsline{toc}{section}{Summary}

Wendelin Werner developed the mathematical theory of Schramm--Loewner evolution (SLE), a family of conformally invariant random curves parametrized by $\kappa > 0$, and used it to compute exact geometric quantities: the Hausdorff dimension $1 + \kappa/8$ of SLE$_\kappa$, the Brownian intersection exponents $(4p^2-1)/24$, and the classification of conformally invariant random sets via the restriction property. The Brownian loop soup and Conformal Loop Ensemble provide the geometric foundations for two-dimensional conformal field theory. SLE has become the central object in the study of scaling limits of critical statistical mechanics models.

\begin{center}\rule{0.5\linewidth}{0.5pt}\end{center}
